\documentclass[10pt,a4paper]{article}
\usepackage[T1]{fontenc}\usepackage[utf8]{inputenc}
\usepackage[margin=17mm]{geometry}
\usepackage{lmodern,microtype,booktabs,xcolor,amsmath}
\usepackage[colorlinks=true,linkcolor=blue,urlcolor=blue]{hyperref}
\setlength{\parskip}{3pt}\setlength{\parindent}{0pt}
\pagestyle{empty}
\newcommand{\CO}{CO\textsubscript{2}}
\begin{document}
{\large\bfseries Supervisor briefing --- routing-policy selection study}\\[2pt]
{\small 21 September 2026 \quad$\cdot$\quad What changed, why, and what we can now claim}

\vspace{4pt}\hrule\vspace{6pt}

\textbf{1. The problem we set out to solve.}
When should an intelligent transportation system use one established routing policy rather than another, and can that choice be learned and evaluated as a \emph{decision} rather than merely predicted?

\textbf{2. What the earlier evidence said.}
The 270-run pilot found dynamic travel-time routing (DTT) best on both criteria in all 30 contexts --- zero adaptive headroom, no study. The 432-run development matrix on a richer 12-signal network was built to create headroom. It did, but the \emph{environmental} part of the design failed: the eco-routing estimator passed aggregate prediction (WAPE 14.71\%) yet failed route ordering (21.62\% selection loss, best route found in 3 of 8 probe panels), and ECO turned out to be outcome-identical to shortest-path routing in 21 of 24 contexts. The earlier conclusion was that the portfolio was insufficient and the study should be withheld.

\textbf{3. What changed, and why.}
That conclusion was too broad. It correctly diagnosed the \emph{environmental} axis as degenerate, but the \emph{policy-selection} problem was intact underneath it. Re-analysing the same 432 runs showed shortest-path routing preferred in all eight low-demand contexts and DTT preferred in all eight high-demand contexts --- a genuine, strong context-to-policy relationship that nobody had looked for, because attention was on the time--\CO\ trade-off that does not exist here. We therefore dropped ECO from the portfolio on evidence, kept it as an audited negative control, and rebuilt the study around demand-regime structure. \textbf{No new simulations were run and no existing result was discarded or reinterpreted.}

\textbf{4. What we can now claim, from completed evidence.}
\begin{itemize}\itemsep1pt
\item \textbf{Policy preference is regime-structured.} The preference between the two main policies reverses with demand, monotonically, in all eight restriction/signal cells.
\item \textbf{The signal regime moves the reversal.} Under balanced signal timing it occurs between 720 and 1200\,veh/h; under arterial-priority timing, between 1200 and 1680\,veh/h. Signal policy, not demand alone, sets where adaptation starts to pay.
\item \textbf{Adaptation is worth something and is recoverable.} Available headroom against the best fixed policy is 0.03018 mean regret units. A depth-3 regression tree, evaluated leave-one-context-out with scales and model refit per fold, recovers \textbf{79.9\%} of it (regret 0.00606) and is optimal in 21 of 24 held-out contexts.
\item \textbf{There is a hard boundary.} Hold out an entire demand level and the selector becomes \emph{worse} than the fixed policy (0.07649 vs 0.03018). The requirement is regime coverage, not model capacity --- a single split on demand already captures 53\%.
\end{itemize}

\textbf{5. The two negative controls we kept.}
ECO was removed from the portfolio because the environment gave it no independent environmental decision dimension --- this is reported, not hidden. The multi-criteria preference layer is nearly inert here: time and \CO\ correlate at 0.982, and across the full weight range from pure-time to pure-emissions \emph{not one} of the 24 selector decisions changes. Both controls are what stop the 79.9\% figure being over-read as a rich multi-objective result.

\textbf{6. What we deliberately did not do.}
We did not repair the emissions observer (not needed once ECO left the portfolio), did not rebuild the network, did not run the planned 2{,}800-run experiment, and did not fit any model more complex than a shallow tree. We also did not claim novelty: a targeted literature audit found demand-dependent routing benefit already established for re-routing decisions, so the paper states differences from named prior studies rather than asserting a first.

\textbf{7. The one open item.}
The demand grid is 480\,veh/h coarse, so the reversal is \emph{bracketed, not localised}. Localising it needs two extra demand levels (960 and 1440\,veh/h) over the existing cells: \textbf{144 runs, about 1.5 minutes of compute}. This was designed and scripted but \textbf{not run}, because the SUMO network and configuration files were not in the material supplied --- a reconstructed network would not be comparable with the 432 completed runs. The paper reports brackets and makes no threshold estimate, so nothing in it depends on this experiment.

\textbf{8. Contribution statement.}
An empirical characterisation of demand- and signal-regime-dependent routing-policy preference in a controlled microsimulation; quantification of the decision headroom that policy switching creates; a held-out, decision-focused evaluation showing an interpretable selector recovers most of that headroom; and identification of an extrapolation boundary beyond which adaptive selection is harmful. No new routing algorithm, no new multi-criteria method, no new learning algorithm.

\vspace{3pt}\hrule\vspace{4pt}
{\small\textbf{Status:} manuscript, figures, tables, bibliography and analysis scripts complete; every reported number regenerates from the exported run table. Outstanding decisions for you: conference template/page limit (none was supplied --- currently generic A4) and the author block.}
\end{document}
